\begin{proposition}[An asymmetric T tree defeats recovery before the next admission]
\label{prop:t-tree-causal-next-admission-stop}
Let $G$ be the six-vertex tree with seed $0$, edge set
\begin{equation}
 \{(0,1),(1,2),(2,3),(2,4),(4,5)\},
 \qquad
 (d_0,d_1,d_2,d_3,d_4,d_5)=(1,2,3,1,2,1),
 \label{eq:t-tree-witness}
\end{equation}
and use the exact transported-center recurrence and complete all-boundary gate
with
\[
 q=\frac15,
 \qquad
 \alpha=\rho=\tau=\frac1{25}.
\]
Every proximal candidate in the resulting execution is strictly positive, so
no displayed event is caused by projection clipping.  The exact chronology is
\begin{center}
\small
\begin{tabular}{@{}ccl@{}}
\toprule
Stage & Action & New vertex, when present \\
\midrule
$1,2$ & admit & $1,2$ \\
$3$ & hold & --- \\
$4$ & admit & $3$ \\
$5$--$8$ & hold & --- \\
$9$ & admit & $4$ \\
$10$--$14$ & hold & --- \\
$15$ & admit & $5$ \\
$16$ & hold & --- \\
$17$ & certify & --- \\
\bottomrule
\end{tabular}
\end{center}
In particular, at stage $4$ the two exposed branch vertices satisfy
\begin{align}
 \frac{\kappa_3}{\sqrt{d_3}}
 &=-\frac{734335307}{179791015625}< -\frac1{625},
 &
 \frac{\kappa_4}{\sqrt{d_4}}
 &=-\frac{223334841}{179791015625}> -\frac1{625},
 \label{eq:t-tree-stage-four-gates}
\end{align}
so only vertex $3$ is admitted.  Vertex $4$ is first admitted at stage $9$,
when
\begin{equation}
 \frac{\kappa_4}{\sqrt{d_4}}
 =-\frac{8110851051223533}{4315546221923828125}
 <-\frac1{625}.
 \label{eq:t-tree-stage-nine-gate}
\end{equation}
The symmetric one-cell Schur load and pivot at the two admissions are
\begin{align}
 (r_4,S_4)&=\left(\frac{167}{23725},\frac{7681}{23725}\right),
 &
 \Delta_4=\frac{r_4^2}{2S_4}
 &=\frac{27889}{364463450},
 \notag\\
 (r_9,S_9)&=\left(\frac{46594}{4800625},\frac{139178}{192025}\right),
 &
 \Delta_9=\frac{r_9^2}{2S_9}
 &=\frac{542750209}{8351767328125}.
 \label{eq:t-tree-schur-drops}
\end{align}

Restart the causal score ledger of
\cref{eq:causal-face-score,eq:causal-recovery-update} at the held stage-$3$
checkpoint with $\mathcal C_3=0$.  It becomes negative after the first
admission's follow-up:
\begin{equation}
 \mathcal C_5
 =-\frac{8829125916972234811}
 {272095189292297363281250}<0.
 \label{eq:t-tree-stage-five-stop}
\end{equation}
That debt does not close before the next admission.  Immediately before its
stage-$9$ Schur credit,
\begin{equation}
 \mathcal C_{9^-}
 =-\frac{68219046926175800933507563}
 {310358575286526679992675781250}<0,
 \label{eq:t-tree-stage-nine-pre-stop}
\end{equation}
and even after the exact drop $\Delta_9$ is realized,
\begin{equation}
 \mathcal C_{9^+}
 =-\frac{33486740471411901479718569}
 {216293380225071907043457031250}<0.
 \label{eq:t-tree-stage-nine-post-stop}
\end{equation}
Every reviewed endpoint from stage $5$ through stage $13$ remains negative;
in particular,
\begin{equation*}
 \mathcal C_{13}
 =-\frac{220753636757299555076490300790138006712609}
 {3633844350039167829883681610226631164550781250}<0.
\end{equation*}
The next held endpoint is the first local GO:
\begin{equation}
 \mathcal C_{14}
 =\frac{1413547795978281227785235568281109142219549159}
 {20440374468970319043095709057524800300598144531250}>0.
 \label{eq:t-tree-stage-fourteen-go}
\end{equation}
Thus stage $13$ is the last local ledger STOP and stage $14$ the first local
ledger GO.  In particular, the assertion that every restarted negative block
recovers before its next safe admission is false for this named exact
recurrence and gate.

The complete finite execution has
\begin{equation}
 J=5,
 \qquad T=17,
 \qquad \nu_{\rm fin}=10,
 \qquad n_{\rm fin}=6,
 \qquad \mathfrak V_T=125.
 \label{eq:t-tree-finite-ledger}
\end{equation}
The scalar audit adds one persistent exact cell and constant control arithmetic
per checkpoint only after the complete gate has materialized $\delta$ and an
admission backend has materialized $\Delta$.  No eleven-coordinate resource
vector is asserted here: the append-only path response and its path-specific
read/write accounting do not automatically extend through the branch in
\cref{eq:t-tree-witness}.
\end{proposition}

\begin{proof}
Use normalized coordinates $z=D^{-1/2}x$.  On every active set $U$, exact
Gaussian elimination solves the restricted stationarity system
\[
 \left(D^{-1/2}QD^{1/2}\right)_{UU}z_U
 =\alpha(\mathbf 1_{\{0\}}-\rho\mathbf 1)_U.
\]
Substitution in the transported-center update, followed by the complete
active and boundary gate, gives the displayed chronology.  Direct exact
evaluation of the two stage-$4$ boundary rows proves
\cref{eq:t-tree-stage-four-gates}; the same calculation at stage $9$ proves
\cref{eq:t-tree-stage-nine-gate}.  The strict positivity of every candidate
checks that these are the literal unprojected recurrence values.

For either one-vertex admission, pad the old restricted optimum by zero and
write the symmetric normalized objective matrix in old/new block form.  Its
new-cell gradient magnitude and Schur pivot are $(r_j,S_j)$, so exact block
elimination gives $\Delta_j=r_j^2/(2S_j)$.  This yields
\cref{eq:t-tree-schur-drops} independently of subtracting the two restricted
objective values.

Here $\Xi_3=0$.  Therefore the causal identity
\cref{eq:causal-recovery-invariant} reads
\[
 \mathcal C_k=\sum_{j\leq k:\,j\text{ admitted}}\Delta_j-\Xi_k,
\]
with the current admission omitted at a pre-reset checkpoint.  Exact squaring
and substitution give
\cref{eq:t-tree-stage-five-stop,eq:t-tree-stage-nine-pre-stop,%
eq:t-tree-stage-nine-post-stop,eq:t-tree-stage-fourteen-go}.  Positive-
denominator comparisons show that every stage-$5$--$13$ balance is negative,
whereas stage $14$ is positive.  Summing the ambient active volumes along the
displayed chronology gives $125$, proving
\cref{eq:t-tree-finite-ledger}.
\end{proof}

\paragraph{Exact scope and computational scaffolding.}
This is a finite exact-real counterexample for the named graph, seed, phase
origin, recurrence, complete gate, and canonical simultaneous boundary-batch
rule.  It refutes only a uniform per-block recovery-before-the-next-admission
claim.  It does not refute the all-history balance, which may contain credits
from stages before the chosen restart; it proves no convergence failure,
objective lower bound, asymptotic chronology or work bound, alternate-order
claim, finite-precision result, or nonpath response complexity.  The exact
checker also exhausts all $3806$ connected labeled rooted graphs on two
through five vertices under this same test and finds no smaller witness.  A
separate reviewed graph-atlas sweep found its first six-vertex witnesses
isomorphic to \cref{eq:t-tree-witness} and edge-minimal among connected
six-vertex witnesses.  Because those finite searches are computational
minimality evidence rather than proof ingredients, none of these minimality
claims is part of the proposition.

The maximizer may keep moving, but its value can be controlled without
locating it.  The next theorem gives the first uniform chronology upper bound
for this exact constant-ratio execution.  It preserves one logarithm and
therefore does not prove the sampled power laws.

\begin{theorem}[Moving-maximum soft chronology upper bound]
\label{thm:path-moving-max-soft-upper}
Under the hypotheses and for the exact execution of
\cref{prop:path-constant-ratio-floor}, define
\begin{equation}
 \mathcal E_{\rm gate}(q)
 :=\frac{q^{12}}{50(1+q^2)^2},
 \qquad
 K_q:=\left\lceil
 \frac{\log\!\left(50(1+q^2)^2q^{-10}\right)}
      {-\log(1-q)}
 \right\rceil.
 \label{eq:path-moving-max-block-length}
\end{equation}
Every visited prefix face undergoes at most $K_q$ consecutive fixed-face
steps before either the complete all-violations gate admits its immediate
successor or the global one-sided certificate holds.  Consequently,
\begin{align}
 T_q
 &\leq (J_q+1)K_q
 \leq\left(\frac5q+1\right)K_q
 =\bigO\!\left(q^{-2}\log\frac1q\right),
 \label{eq:path-moving-max-stage-upper}\\
 \mathfrak V_{T_q}
 &\leq\frac5qT_q
 =\bigO\!\left(q^{-3}\log\frac1q\right).
 \label{eq:path-moving-max-volume-upper}
\end{align}
For the named literal full-prefix exact-cell implementation with one terminal
external certificate/output return, the resulting eleven-coordinate upper
ledger is
\begin{equation}
 \begin{aligned}
 (&C_{\rm adj},R_{\rm adj},R_{\rm int},C_{\rm pre},C_{\rm ctl},C_{\rm rec},
 C_{\rm resp},M_{\rm pers},M_{\rm tmp},C_{\rm mat},C_{\rm emit})\\
 =\bigl(&\Theta(q^{-1}),\Theta(q^{-1}),1,0,
 O(q^{-3}\log(1/q)),O(q^{-3}\log(1/q)),
 O(q^{-3}\log(1/q)),\\
 &\Theta(q^{-1}),O(q^{-1}),O(q^{-3}\log(1/q)),
 \Theta(q^{-1})\bigr).
 \end{aligned}
 \label{eq:path-moving-max-eleven-upper}
\end{equation}
The bounds hold regardless of how often or where the correction maximizer
switches.  They do not prove $T_q=\Theta(q^{-2})$ or
$\mathfrak V_{T_q}=\Theta(q^{-3})$.  Removal of the logarithm, whether by a
uniform fixed-face argument or by an aggregate potential, and matching global
lower powers remain open.
\end{theorem}

\begin{proof}
Write $x_{U}^{\star}$ for the interior restricted optimum on a visited path
prefix $U$.  Interiority holds at the seed face and is preserved by every
strong path admission, as in the proof of
\cref{prop:path-implicit-transport-ledger}.  At every post-transport state,
the weighted Schur ledger and its unweighted telescope imply
\begin{align}
 \mathcal E_U(x,v)
 &\leq \mathcal E_{U_0}(0,0)
       +G_\rho(x_{U_0}^{\star})-G_\rho(x_U^{\star})\notag\\
 &=G_\rho(0)-G_\rho(x_U^{\star})
   +\frac{\alpha}{2}\norm{x_{U_0}^{\star}}_2^2
 \leq q^2.
 \label{eq:path-moving-max-energy-cap}
\end{align}
Indeed, on this endpoint-seeded path the positive part of $c_\rho$ consists
only of the seed entry and has Euclidean norm at most $\alpha$.  Since
$Q\succeq\alpha I$ and feasible vectors are nonnegative,
\begin{equation*}
 G_\rho(0)-G_\rho(x_\rho^\star)
 \leq\frac{\norm{(c_\rho)_+}_2^2}{2\alpha}
 \leq\frac{q^2}{2}.
\end{equation*}
Moreover
$x_{U_0}^{\star}=\alpha(1-\rho)/a$ with $a=(1+q^2)/2$, so its center term is
at most $2q^6\leq q^2/2$ for $q\leq1/4$.  Restricted optimality and
$G_\rho(x_U^\star)\geq G_\rho(x_\rho^\star)$ prove the last inequality in
\cref{eq:path-moving-max-energy-cap}.

Now put
\begin{equation*}
 p:=D_U^{-1/2}x,
 \qquad p^\star:=D_U^{-1/2}x_U^\star,
 \qquad e:=\norm{p-p^\star}_\infty.
\end{equation*}
The normalized active KKT operator is
\begin{equation*}
 \widehat Q_U:=D_U^{-1/2}Q_{UU}D_U^{1/2}.
\end{equation*}
Every row of $\widehat Q_U$ has diagonal magnitude
$a=(1+q^2)/2$ and off-diagonal absolute row sum at most
$\eta=(1-q^2)/2$.  Hence
$\norm{\widehat Q_U}_{\infty\to\infty}\leq a+\eta=1$.  Because
$\kappa_{\rho,U}(x_U^\star)=0$, the moving correction obeys the pointwise
bound
\begin{equation}
 \delta(x)
 \leq\frac{e}{q^2}.
 \label{eq:path-moving-max-delta-energy-a}
\end{equation}
The normalized safe envelope is $[p-\delta(x)\one]_+$, and $p^\star\geq0$;
therefore
\begin{equation}
 \norm{D_U^{-1/2}(\ell-x_U^\star)}_\infty
 \leq (1+q^{-2})e.
 \label{eq:path-moving-max-envelope-error}
\end{equation}
Finally, the quadratic gap in $\mathcal E_U$ is at least
$\alpha\norm{x-x_U^\star}_2^2/2$.  Since every path degree is at least one,
\begin{equation}
 e\leq\frac{\sqrt{2\mathcal E_U(x,v)}}q.
 \label{eq:path-moving-max-delta-energy-b}
\end{equation}

If $\mathcal E_U(x,v)\leq\mathcal E_{\rm gate}(q)$, then
\cref{eq:path-moving-max-envelope-error,eq:path-moving-max-delta-energy-b}
give
\begin{equation*}
 \norm{D_U^{-1/2}(\ell-x_U^\star)}_\infty
 \leq\frac{q^3}{5}=\alpha\tau.
\end{equation*}
Applying the same unit row-sum bound to the active KKT vector shows that every
active row satisfies
$\kappa_{\rho,i}(\ell)/\sqrt{d_i}\geq-\alpha\tau$.
On a path prefix, every exterior row beyond the immediate successor has
normalized KKT value $\alpha\rho>0$.  The immediate successor is therefore
either a strong violation and is admitted by the complete gate, or it too
satisfies the one-sided threshold, in which case
\cref{eq:global-certificate} holds.  This conclusion uses the maximum over
the whole old face through \cref{eq:path-moving-max-delta-energy-a}; it never
assumes that its maximizer is the frontier or the seed.

During $k$ consecutive steps on one fixed face,
\cref{lem:fixed-face-estimate-center} and
\cref{eq:path-moving-max-energy-cap} give
\begin{equation*}
 \mathcal E_U(x_k,v_k)\leq(1-q)^kq^2.
\end{equation*}
The definition of $K_q$ makes the right-hand side at most
$\mathcal E_{\rm gate}(q)$.  Hence a face cannot survive more than $K_q$
steps without admission or termination.  There are $J_q+1$ visited faces and
$J_q\leq1/\rho=5/q$, proving
\cref{eq:path-moving-max-stage-upper}.  Safe support also gives
$\vol(U_t)\leq1/\rho=5/q$ at every stage, which proves
\cref{eq:path-moving-max-volume-upper}.  Since
$-\log(1-q)\geq q$, one has
$K_q=\bigO(q^{-1}\log(1/q))$ uniformly for $0<q\leq1/4$.
Substitution into \cref{eq:path-constant-ratio-eleven-vector}, together with
$J_q,\nu_{\rm fin}=\Theta(q^{-1})$, proves
\cref{eq:path-moving-max-eleven-upper}.
\end{proof}

\paragraph{Exact scope of the upper bound.}
The theorem concerns only the named zero-start, internally gated,
transported-center exact-real execution on finite endpoint paths and its
literal full-prefix implementation.  It retains one terminal external return
and the note-scoped PPR error $\rho+\tau=2q/5$ from
\cref{prop:path-constant-ratio-floor}.  It neither proves a lower bound at
the displayed upper powers nor applies to the literal zero-padded recurrence,
branching or cyclic graphs, intermediate external emissions, finite-precision
arithmetic, numerical stability, or bit complexity.

\paragraph{Audited terminal-face candidate: exact roots and finite scaffolding.}
Fix an integer $m\geq2$ and put
\begin{equation}
 q_m:=\frac1{16m},
 \qquad
 \alpha_m:=q_m^2,
 \qquad
 \rho_m=\tau_m:=\frac{q_m}{5}.
 \label{eq:path-terminal-candidate-parameters}
\end{equation}
The candidate uses the zero-start, internally gated, transported-center
algorithm and parameter ratio studied in \cref{sec:path-constant-ratio}, but
runs it on the complete endpoint-seeded path
$P_m=(v_0,\ldots,v_m)$, with ambient degrees, the complete all-violations
gate, and one terminal external certificate/output return.  It is not an
instance of \cref{prop:path-constant-ratio-floor}, whose path-length premise
is longer; the finite trajectory is replayed directly below.
The open question is whether its final full-path face can require
$\Omega(q_m^{-1}\log(1/q_m))$ consecutive steps.  The following conditional
calculation is rigorous, but it does not answer that question.

\begin{lemma}[Conditional characteristic roots on a full path face]
\label{lem:path-full-face-linearized-roots}
Let the current face be all of $P_m$, and consider consecutive fixed-face
steps of \cref{eq:center-prox-step,eq:center-update} on which the
nonnegativity projection is inactive.  Set
\begin{equation*}
 \begin{aligned}
 a&:=\frac{1+q^2}{2},
 &\eta&:=\frac{1-q^2}{2},
 &P_{\mathrm{row},m}&:=D^{-1}A=P^\top,\\
 \widehat Q&:=D^{-1/2}QD^{1/2}
             =aI-\eta P_{\mathrm{row},m},
 &B&:=I-\widehat Q.
 \end{aligned}
\end{equation*}
For the affine normalized residual states
\[
 r_k^x:=D^{-1/2}(Qx_k-c_\rho),
 \qquad
 r_k^v:=D^{-1/2}(Qv_k-c_\rho),
\]
one fixed-face step obeys
\begin{equation}
 r_{k+1}^x
 =\frac{B}{1+q}(r_k^x+q r_k^v),
 \qquad
 q r_{k+1}^v=r_{k+1}^x-(1-q)r_k^x.
 \label{eq:path-full-face-two-state}
\end{equation}
Consequently, for $k\geq1$,
\begin{equation}
 r_{k+1}^x
 =B\left(
   \frac{2}{1+q}r_k^x
   -\frac{1-q}{1+q}r_{k-1}^x
 \right).
 \label{eq:path-full-face-residual-recurrence}
\end{equation}
Here $P=AD^{-1}$ is the shared column-walk convention.  The eigenvalues of
$P_{\mathrm{row},m}$ are
$\mu_h=\cos(h\pi/m)$, $0\leq h\leq m$.  On the $h$th cosine mode, the two
characteristic roots of \cref{eq:path-full-face-residual-recurrence} are
\begin{equation}
 \zeta_h^\pm
 =(1-q)\cos\!\left(\frac{h\pi}{2m}\right)
  \exp\!\left(\pm\frac{\mathrm i h\pi}{2m}\right).
 \label{eq:path-full-face-linearized-roots}
\end{equation}
\end{lemma}

\begin{proof}
Inactivity of the projection gives
$x_{k+1}=y_k-(Qy_k-c_\rho)$.  Since
$y_k=(x_k+qv_k)/(1+q)$ and the affine coefficients sum to one,
normalization yields the first identity in
\cref{eq:path-full-face-two-state}.  Multiplying the center update by $q$
gives $qv_{k+1}=x_{k+1}-(1-q)x_k$, which yields its second identity.
Eliminating $r_k^v$ proves
\cref{eq:path-full-face-residual-recurrence}.

For $\theta_h=h\pi/m$, the $h$th eigenvalue of $B$ is
\[
 b_h=\eta(1+\cos\theta_h)
     =(1-q^2)\cos^2(\theta_h/2).
\]
The scalar characteristic polynomial is therefore
\[
 z^2-2(1-q)\cos^2(\theta_h/2)z
 +(1-q)^2\cos^2(\theta_h/2),
\]
whose roots are exactly
\cref{eq:path-full-face-linearized-roots}.
\end{proof}

\paragraph{Finite exact scaffolding.}
The direction-local script \texttt{check\_round013.py} uses rational
arithmetic for
\[
 (m,q_m)\in\{(2,1/32),(4,1/64),(8,1/128)\}.
\]
For those three individual instances it replays the actual zero-start
complete all-violations gate, verifies respectively $m$ successive singleton
admissions, checks \cref{eq:path-full-face-two-state} exactly on the terminal
face, and verifies inactive projection, an unclipped moving global envelope,
and non-certification for respectively $16$, $32$, and $64$ terminal steps.
The checker inspects these internal prefixes and makes no intermediate
external return.  These are finite exact witnesses only: the tested terminal
length is $1/(2q_m)$, not the conjectured
$\Theta(q_m^{-1}\log(1/q_m))$ range, and three sizes supply no asymptotic
inequality.

\paragraph{Failed lower-bound audit and rigorous GO/STOP.}
An earlier draft promoted the candidate to a logarithmic lower theorem.  Two
independent audits rejected that promotion.  The characteristic roots specify
conditional propagation but not the cosine coefficients of the actual state
upon entry to the full face.  No uniform lower bound on a surviving mode or
anti-cancellation estimate for the moving residual range was derived.
In particular, deleting nonpositive packet fragments is not monotone for
oscillation.  The proposed rapid-front Pascal induction also omitted explicit
base and endpoint-row derivations, and neither projection inactivity nor the
unclipped-envelope margin was proved uniformly through
$k_0=\Theta(q_m^{-1}\log(1/q_m))$.

Thus there is no proved terminal logarithmic block, no asymptotic lower
eleven-vector, and no refutation of a uniform
$K_{\rm face}=O(q^{-1})$ bound.  The existing upper ledger
\cref{eq:path-moving-max-eleven-upper} remains the only rigorous
constant-ratio chronology ledger.  A valid lower-bound GO must derive the
literal full-face entry coefficients, prove the actual zero-start
all-violations trajectory remains in the conditional regime through $k_0$,
and establish a moving-global-residual range lower bound with explicit
anti-cancellation.  The alternative GO is an admission-specific energy,
normalized-error, or space--time potential proving
$K_{\rm face}=O(q^{-1})$, $T=O(q^{-2})$, and
$\mathfrak V_T=O(q^{-3})$.  Either route must retain the full moving global
correction; a frontier or seed surrogate is not admissible.

The exact STOP in
\cref{prop:path-monotone-correction-potential-fails} additionally rules out
treating the correction itself as monotone held-face debt or replacing the
reviewed pointwise factor $q^{-2}$ by the displayed coefficient-one
$q^{-1}$ factor.  It does not rule out a nonmonotone or phase-aware aggregate
argument.  The next explicit attempt
$\Phi_t^{\rm bank}=\delta_t+q^{-1}e_t$ does absorb the stage-$6$--$7$ spike,
but \cref{prop:path-correction-error-bank-potential-fails} refutes its
held-pair monotonicity at stages $9$--$10$.  That finite STOP leaves longer
blocks and cross-face amortization open.  Moreover,
\cref{prop:path-no-constant-coefficient-bank} derives the exact local
coefficient constraints from both pairs and proves that no fixed
$c\geq0$ makes $\delta_t+cq^{-1}e_t$ nonincreasing on both.  This pairwise
incompatibility does not define a global block horizon or exclude a
time-varying, face-dependent, or signed-phase rule.  The canonical endpoint
test in \cref{prop:path-tight-pair-schur-bank} gives one finite positive
reset identity: after squaring the bank, the exact stage-$8$ coefficient
switch and optimum reset are paid by the restricted-optimum drop.  Its raw
linear version fails the same unit-shock comparison.  This keeps held-face
evolution and the frozen-candidate reset separate.  The extension in
\cref{prop:path-tight-pair-local-chain-stop} now evaluates both adjacent
fixed-face steps: the squared-bank-plus-reserve account already rises on the
candidate-producing step $7\to8^-$, falls across the paid reset, rises again
from $8^+$ to $9$, and has positive net stage-$7$--$9$ change.  Both
adjacent-step failures persist throughout the held-pair-feasible coefficient
rectangle.  This closes that pointwise local-chain attempt, not longer
aggregate blocks or a multi-face telescope.  The first tested longer account
in \cref{prop:path-tight-pair-recovery-block} is positive but finite: it
accumulates only the subsequent fixed-$U_4$ squared-bank decreases.  The
reserve is insufficient through stage $12$ and first pays at stage $13$,
giving a rectangle-uniform $7\to13$ endpoint decrease.  This stage-$7$--$13$
single-admission block neither supplies a face-general closing time nor
iterates across admissions.
The causal condition in
\cref{prop:path-causal-two-admission-recovery} supplies the next exact
interface without claiming such a time bound.  Its score $\Xi=\delta^2$ is
observable on every face, and its balance credits only already computed score
decreases and already realized restricted-optimum drops while debiting every
increase.  On the same trace, unused stage-$4$ credit remains solvent through
the next canonical admission at stage $8$; without that cross-state carry, the
stage-$8$ account first recovers at stage $12$.  This is a finite conditional
scalar-ledger GO and an exact cross-state cancellation, not a convergence or
work certificate and not a proof that another graph,
face, order, or admission sequence has nonnegative balance.  The asymmetric
T-tree calculation in
\cref{prop:t-tree-causal-next-admission-stop} closes precisely that next
falsification target: after restarting at stage $3$, its debt is still
negative before and after the next admission.  It retires uniform
recovery-before-the-next-admission for this recurrence, while leaving an
all-history ledger and graph-structural solvency conditions as the next
targets.  The former target is refuted by the zero-start three-admission
execution in \cref{prop:three-admission-all-history-stop}; the latter remains
open.

\paragraph{Numerical scaffolding, not a theorem.}
An independent floating-point implementation of the same recurrence returned
\begin{center}
\begin{tabular}{@{}c|rrrr@{}}
$q$ & $J$ & $T$ & $\mathfrak V_T$ &
$(qJ,q^2T,q^3\mathfrak V_T)$\\
\midrule
$1/8$  & $7$  & $39$   & $423$    & $(0.875,0.609,0.826)$\\
$1/16$ & $16$ & $168$  & $3978$   & $(1.000,0.656,0.971)$\\
$1/32$ & $34$ & $772$  & $42106$  & $(1.063,0.754,1.285)$\\
$1/64$ & $71$ & $3257$ & $363001$ & $(1.109,0.795,1.385)$
\end{tabular}
\end{center}
These values motivate the open matching global lower bounds and a possible
aggregate removal of the logarithm in
\cref{thm:path-moving-max-soft-upper}.  They are not used in any theorem and
do not prove that the sampled constants or exact power laws persist as
$q\downarrow0$.  The audited short-path candidate above likewise supplies no
asymptotic terminal-face lower bound.

\begin{conjecture}[One-sided expanding-subspace acceleration]
\label{conj:one-sided-continuation}
For the recurrence \cref{eq:dynamic-y,eq:dynamic-step,eq:dynamic-gate}, or for
a closely related projected Chebyshev recurrence, universal constants $C,c>0$
exist such that
\begin{equation}
 \norm{D^{-1/2}(x_\rho^\star-\ell^{(t)})}_\infty
 \leq C\exp(-c\sqrt\alpha\,t)+\tau.
 \label{eq:desired-continuation}
\end{equation}
\end{conjecture}

If \cref{conj:one-sided-continuation} holds, then
\begin{equation}
 t=\bigO\!\left(\frac1{\sqrt\alpha}\log\frac1\tau\right),
 \qquad
 \mathcal W
 =\bigO\!\left(
   \frac1{\rho\sqrt\alpha}\log\frac1\tau
 \right).
 \label{eq:conditional-final}
\end{equation}
With $\rho=\tau=\epsppr/2$, this is the sought
$\softO(1/(\sqrt\alpha\epsppr))$ bound.

\paragraph{Why existing analyses do not immediately prove it.}
The nonhomogeneous \LOCCH{} formula of \citet{zhou2024} controls a full signed
inactive residual.  The safe gate controls only the negative part of the KKT
residual at a corrected lower envelope.  Outside the optimal support, positive
KKT slack can be dense and is harmless for the obstacle problem, but it enters
a two-sided residual norm.  Energy contraction alone also does not imply the
uniform active-volume invariant.  Finally, \cref{prop:shock} supplies an
upper perturbation bound only with a possible $1/\alpha$ loss.

\paragraph{A weaker sufficient target.}
It is enough to strengthen \cref{thm:amortized} so that discrete expansion
overhead is charged to newly admitted volume:
\begin{equation}
  \mathcal W_{\mathrm{exp}}
  \leq C\sum_t\vol(U_{t+1}\setminus U_t)
  \leq\frac{C}{\rho},
  \label{eq:new-volume-charge}
\end{equation}
rather than $N_{\mathrm{exp}}/\rho$.  Combining
\cref{eq:new-volume-charge} with the continuous telescoping term already gives
the target up to a logarithm once the number and placement of recurrence
sweeps are controlled.  Proposition
\cref{prop:path-implicit-transport-ledger} proves this charge for the
incremental transport overhead on endpoint paths, but leaves every old-face
recurrence/response sweep in $\mathfrak V_T$ and does not bound $T$.
